\documentclass[12pt,a4paper]{article}
% Drake_preamble.tex - LaTeX Preamble Template
% 作者: 謝慕揚, MD, PhD, FESC
% 
% 使用說明:
% 1. 表格請使用 longtable，不要有垂直分隔線 |
% 2. 表格內換行使用 \newline，不要使用 \\
% 3. 藥物名稱、檢驗、檢查請維持英文
% 4. Reference 格式請使用 \href 加上驗證過的 DOI

% Including essential packages for Chinese typesetting
\usepackage{xeCJK}
\usepackage{fontspec}
% Configuring Chinese font with Noto Serif CJK TC, fallback to SimSun
\IfFontExistsTF{Noto Serif CJK TC}{
  \setCJKmainfont{Noto Serif CJK TC}
  \setCJKsansfont{Noto Serif CJK TC}
  \setCJKmonofont{Noto Serif CJK TC}
}{\setCJKmainfont{SimSun}
  \setCJKsansfont{SimSun}
  \setCJKmonofont{SimSun}
}

% Including other necessary packages
\usepackage{geometry}
\usepackage{titlesec}
\usepackage{enumitem}
\usepackage{xcolor}
\usepackage{colortbl}
\usepackage{tcolorbox}
\usepackage{booktabs}
\usepackage{longtable}
\usepackage{array}
\usepackage{graphicx}
\usepackage{tikz}
\usepackage{fancyhdr}
\usepackage{multicol}
\usepackage{datetime2}
\usepackage{hyperref}
\usepackage{mdframed}
\usepackage{amsmath}
\usepackage{amssymb}
\usepackage{multirow}

\hypersetup{
    colorlinks=true,
    linkcolor=emergency_blue,
    citecolor=emergency_blue,
    urlcolor=emergency_blue,
    pdfborder={0 0 0}
}

% Defining page geometry
\geometry{
  top=2.5cm,
  bottom=2.5cm,
  left=2cm,
  right=2cm,
  headheight=25pt
}

% Adjusting topmargin to compensate for increased headheight
\addtolength{\topmargin}{-10pt}

% Defining colors
\definecolor{emergency_red}{RGB}{186,24,27}
\definecolor{emergency_blue}{RGB}{0,114,188}
\definecolor{emergency_gray}{RGB}{120,120,120}
\definecolor{highlight_yellow}{RGB}{255,250,205}

% Setting title formats
\titleformat{\section}[block]{\Large\bfseries\color{emergency_red}}{\thesection}{10pt}{}
\titleformat{\subsection}[block]{\large\bfseries\color{emergency_blue}}{\thesubsection}{8pt}{}
\titleformat{\subsubsection}[block]{\normalsize\bfseries\color{emergency_gray}}{\thesubsubsection}{6pt}{}

% Defining custom environment for important notes
\newmdenv[
    linecolor=emergency_red,
    backgroundcolor=highlight_yellow,
    linewidth=2pt,
    topline=true,
    bottomline=true,
    leftline=true,
    rightline=true,
    innertopmargin=10pt,
    innerbottommargin=10pt,
    innerleftmargin=10pt,
    innerrightmargin=10pt
]{importantbox}

% Defining note command with tcolorbox
\newcommand{\note}[1]{
    \par
    \noindent
    \begin{tcolorbox}[
        colback=emergency_blue!5!white,
        colframe=emergency_blue,
        title=\textbf{重點提醒},
        fonttitle=\bfseries,
        boxrule=0.5pt,
        arc=3pt,
        left=6pt,
        right=6pt,
        top=6pt,
        bottom=6pt
    ]
    #1
    \end{tcolorbox}
}

% Key findings box
\newcommand{\keyfinding}[1]{
    \par
    \noindent
    \begin{tcolorbox}[
        colback=yellow!10!white,
        colframe=orange!75!black,
        title=\textbf{核心發現摘要},
        fonttitle=\bfseries,
        boxrule=0.5pt,
        arc=3pt
    ]
    #1
    \end{tcolorbox}
}

% Defining custom date format for ISO 8601
\DTMnewdatestyle{iso}{
  \renewcommand{\DTMdisplaydate}[4]{\number##1-\number##2-\number##3}
  \renewcommand{\DTMDisplaydate}[4]{\number##1-\number##2-\number##3}
}
\DTMsetdatestyle{iso}
\newcommand{\mydate}{\DTMtoday}


% Custom header for this document
\pagestyle{fancy}
\fancyhf{}
\fancyhead[L]{\textcolor{emergency_red}{Intensive Care Medicine 教學講義}}
\fancyhead[R]{\textcolor{emergency_blue}{AKI \& RRT Timing}}
\fancyfoot[C]{\textcolor{emergency_gray}{\thepage}}
\renewcommand{\headrulewidth}{1pt}
\renewcommand{\headrule}{\hbox to\headwidth{\color{emergency_red}\leaders\hrule height \headrulewidth\hfill}}

\begin{document}

\begin{center}
{\LARGE\bfseries\textcolor{emergency_red}{Intensive Care Medicine}}\\[0.5cm]
{\Large\textbf{Acute Kidney Injury: When and How to Start Renal Replacement Therapy}}\\[0.3cm]
{\large 急性腎損傷：何時及如何開始腎臟替代療法}\\[0.5cm]
{\normalsize Intensive Care Med 2025;51:1172-1175. DOI: 10.1007/s00134-025-07933-x}\\[0.3cm]
{\normalsize 整理：謝慕揚 MD, PhD, FESC \quad 日期：\mydate}
\end{center}

\keyfinding{
\textbf{核心訊息：}對於 KDIGO stage 2-3 AKI 但無 life-threatening complications 的病人，delayed RRT strategy 可使高達 50\% 病人避免透析，且不影響存活率。

\textbf{臨床建議：}
\begin{itemize}
    \item 無緊急適應症時，採取 watchful waiting 策略
    \item 但對於 persistent AKI（oliguria 或 BUN >112 mg/dL），不應延遲超過 Stage 3 AKI 發生後 72 小時
    \item CRRT 建議劑量 20-25 mL/kg/h，優先使用 regional citrate anticoagulation
\end{itemize}
}

\tableofcontents
\newpage

%==============================================================
\section{文獻概述}
%==============================================================

本篇為 Intensive Care Medicine「What's New in Intensive Care」系列文章，由法國 Nimes 大學醫院 Barbar 等人撰寫，綜合最新研究證據，針對 ICU 病人 AKI 的 RRT 啟動時機與方式提供實用建議。

\subsection{作者資訊}
\textbf{通訊作者：}Saber Davide Barbar, MD\\
\textbf{機構：}Division of Anesthesia Critical Care, Pain and Emergency Medicine, Nimes University Hospital, France

%==============================================================
\section{背景}
%==============================================================

\begin{itemize}
    \item AKI 影響高達 \textbf{50\%} 的 ICU 病人
    \item 與較高的死亡率及 progressive kidney disease 風險相關
    \item RRT 啟動時機及方式仍有爭議
\end{itemize}

%==============================================================
\section{何時開始 RRT}
%==============================================================

\subsection{無 Life-threatening Complications 的 AKI}

\note{
\textbf{關鍵證據：}多個 RCT（AKIKI, IDEAL-ICU, STARRT-AKI, ELAIN）的核心訊息：

在無 life-threatening complication 的情況下，即使 KDIGO stage 2-3 AKI，early RRT 並不能改善存活率。
}

\subsubsection{Delayed Strategy 的優勢}

\begin{itemize}
    \item 高達 \textbf{50\%} 的病人可避免 RRT（因腎功能自發性改善）
    \item 各研究中 delayed arm 的病人，RRT 啟動時間為 enrollment 後 \textbf{31-57 小時}
\end{itemize}

\subsubsection{Early RRT 的潛在危害}

\begin{itemize}
    \item 在無 oligo-anuria 的病人，early RRT 可能導致 \textbf{excess mortality}
    \item 可能增加 \textbf{long-term dialysis dependence} 風險，特別是有 pre-existing CKD 的病人
    \item 機轉：RRT-associated hemodynamic instability 可能損害腎臟修復及內源性腎功能恢復
\end{itemize}

\subsubsection{AKIKI-2 Trial 的重要發現}

\begin{longtable}{p{4cm}p{10cm}}
\toprule
\textbf{比較} & \textbf{結果} \\
\midrule
\endfirsthead

\toprule
\textbf{比較} & \textbf{結果} \\
\midrule
\endhead

\bottomrule
\endfoot

Delayed vs More Delayed & "More delayed" strategy 並未帶來更多 RRT-free days \\
\midrule
死亡率 & "More delayed" 有較高死亡率的趨勢 \\
\midrule
Comatose 病人 & More delayed RRT 啟動與較低的 coma to awakening 轉換機率相關 \\
\bottomrule
\end{longtable}

\subsection{Life-threatening Complications：緊急 RRT 適應症}

\begin{longtable}{p{5cm}p{9cm}}
\toprule
\textbf{緊急適應症} & \textbf{處理建議} \\
\midrule
\endfirsthead

\toprule
\textbf{緊急適應症} & \textbf{處理建議} \\
\midrule
\endhead

\bottomrule
\endfoot

\textbf{Severe Hyperkalemia} &
K+ 5.5-6.5 mmol/L + ECG changes + refractory to medical treatment\newline
可先嘗試 diuretics（若有 volume overload）或 novel potassium binders \\
\midrule

\textbf{Metabolic Acidosis} &
pH ≤7.20 + KDIGO stage 2-3 AKI\newline
先給 sodium bicarbonate（250 mL of 4.2\%）\newline
\textbf{若 pH 仍 <7.20，不應延遲 RRT} \\
\midrule

\textbf{Fluid Overload} &
最難客觀定義的適應症\newline
Physical exam、fluid balance data、chest X-ray、oxygenation status 整合判斷\newline
可輔以 bedside ultrasound 及 bioimpedance \\
\bottomrule
\end{longtable}

\note{
\textbf{Fluid Overload 定義建議：}
\begin{itemize}
    \item Worsening pulmonary oedema
    \item PaO2/FiO2 <200 mmHg
    \item Fluid balance >10\% of body weight
    \item 對 diuretics 無反應
\end{itemize}
}

%==============================================================
\section{如何開始 RRT}
%==============================================================

\subsection{RRT Modality 選擇：CRRT vs IHD}

\begin{longtable}{p{3cm}p{5.5cm}p{5.5cm}}
\toprule
\textbf{項目} & \textbf{CRRT} & \textbf{IHD} \\
\midrule
\endfirsthead

\toprule
\textbf{項目} & \textbf{CRRT} & \textbf{IHD} \\
\midrule
\endhead

\bottomrule
\endfoot

適用對象 & 使用 inotropes/vasopressors\newline Significant volume overload & Less severe disease\newline SOFA score <10 \\
\midrule
存活率 & 大型 RCT 未顯示差異 & 大型 RCT 未顯示差異 \\
\midrule
腎臟預後 & STARRT-AKI 次級分析：\newline 較低 90-day RRT dependence & AKIKI/IDEAL-ICU 次級分析：\newline SOFA <10 時有較好 60-day survival \\
\bottomrule
\end{longtable}

\subsection{CRRT 處方建議}

\begin{longtable}{p{5cm}p{9cm}}
\toprule
\textbf{參數} & \textbf{建議} \\
\midrule
\endfirsthead

\toprule
\textbf{參數} & \textbf{建議} \\
\midrule
\endhead

\bottomrule
\endfoot

\textbf{Small molecule clearance} & 20-25 mL/kg/h（標準劑量）\newline 較低劑量是否同樣有效仍待研究 \\
\midrule

\textbf{Anticoagulation} & \textbf{Regional citrate > Heparin}\newline 較長 filter life、較低出血風險 \\
\midrule

\textbf{Hemofiltration vs Hemodialysis} & 尚無定論\newline Hemofiltration 透過 convection 移除較大分子\newline Hemodialysis 依賴 diffusive clearance \\
\midrule

\textbf{Ultrafiltration rate} & 仍有爭議\newline 積極移水可較快達到 euvolemia，但代價是 hemodynamic instability \\
\bottomrule
\end{longtable}

%==============================================================
\section{實用決策流程}
%==============================================================

\begin{center}
\fbox{\parbox{0.9\textwidth}{
\textbf{AKI KDIGO ≥2 in ICU ± MOF}

$\downarrow$

\textbf{Step 1：評估是否有 Life-threatening / Emergency criteria for RRT}

\begin{itemize}
    \item Hyperkalemia >6.5 + ECG signs + refractory to medical treatment
    \item Metabolic acidosis (pH ≤7.15) unresponsive to 100 mL of 4.2\% bicarb
    \item Fluid overload* + refractory to diuretics
\end{itemize}

$\downarrow$ Yes → \textbf{Immediate RRT}

$\downarrow$ No → \textbf{Dynamic reassessment every 8h}

$\downarrow$

\textbf{Step 2：若無緊急適應症，等待但不超過 72 小時}

若持續 AKI（oliguria 或 BUN >112 mg/dL）→ \textbf{Delayed RRT after maximum 72h}

\vspace{0.3cm}
\small *Fluid overload: Worsening pulmonary oedema, PaO2/FiO2 <200 mmHg, or fluid balance >10\% body weight
}}
\end{center}

%==============================================================
\section{重要臨床試驗摘要}
%==============================================================

\begin{longtable}{p{3cm}p{4cm}p{7cm}}
\toprule
\textbf{Trial} & \textbf{比較} & \textbf{主要發現} \\
\midrule
\endfirsthead

\toprule
\textbf{Trial} & \textbf{比較} & \textbf{主要發現} \\
\midrule
\endhead

\bottomrule
\endfoot

AKIKI (2016) & Early vs Delayed RRT & Delayed 可避免 49\% RRT，存活率無差異 \\
\midrule
IDEAL-ICU (2018) & Early vs Delayed RRT in sepsis & 無存活差異，early futility stop \\
\midrule
STARRT-AKI (2020) & Accelerated vs Standard & 無存活差異，accelerated 有更多 dialysis dependence \\
\midrule
ELAIN (2016) & Early vs Delayed & 唯一顯示 early 優勢的試驗，但多數病人有 volume overload \\
\midrule
AKIKI-2 (2021) & Delayed vs More delayed & More delayed 無額外好處，可能有害 \\
\midrule
BICAR-ICU (2018) & Sodium bicarbonate for metabolic acidosis & AKI + pH ≤7.20：bicarb 可降低死亡率及 RRT 需求 \\
\bottomrule
\end{longtable}

%==============================================================
\section{結論與建議}
%==============================================================

\subsection{When to Initiate RRT}

\begin{enumerate}
    \item \textcolor{red}{\textbf{無緊急適應症}}：採取 watchful waiting + RRT deferral 策略
    \item \textcolor{red}{\textbf{存活率}}：RRT 啟動時機可能與存活率無關，但 delayed 可讓部分病人避免 RRT
    \item \textcolor{red}{\textbf{不應過度延遲}}：Unresolving AKI（persistent oliguria 或 BUN >112 mg/dL）不應等待超過 Stage 3 AKI 發生後 72 小時
\end{enumerate}

\subsection{How to Initiate RRT}

\begin{enumerate}
    \item \textcolor{blue}{\textbf{Modality 選擇}}：根據血流動力學狀態決定
    \begin{itemize}
        \item 使用 inotropes/vasopressors 或嚴重 volume overload → CRRT
        \item Less severe disease (SOFA <10) → IHD 可能是安全替代
    \end{itemize}
    \item \textcolor{blue}{\textbf{CRRT 劑量}}：20-25 mL/kg/h
    \item \textcolor{blue}{\textbf{Anticoagulation}}：Regional citrate 優於 heparin
\end{enumerate}

%==============================================================
\section{縮寫對照表}
%==============================================================

\begin{longtable}{p{3cm}p{11cm}}
\toprule
\textbf{縮寫} & \textbf{全名} \\
\midrule
\endfirsthead

\toprule
\textbf{縮寫} & \textbf{全名} \\
\midrule
\endhead

\bottomrule
\endfoot

AKI & Acute Kidney Injury (急性腎損傷) \\
RRT & Renal Replacement Therapy (腎臟替代療法) \\
KDIGO & Kidney Disease: Improving Global Outcomes \\
CRRT & Continuous Renal Replacement Therapy (連續性腎臟替代療法) \\
IHD & Intermittent Hemodialysis (間歇性血液透析) \\
CKD & Chronic Kidney Disease (慢性腎臟病) \\
BUN & Blood Urea Nitrogen (血中尿素氮) \\
SOFA & Sequential Organ Failure Assessment \\
MOF & Multiorgan Failure (多重器官衰竭) \\
\bottomrule
\end{longtable}

%==============================================================
\section{參考文獻}
%==============================================================

\begin{enumerate}
    \item Barbar SD, Wald R, Quenot JP. Acute kidney injury: when and how to start renal replacement therapy. \href{https://doi.org/10.1007/s00134-025-07933-x}{\textit{Intensive Care Med}. 2025;51:1172-1175.}

    \item Gaudry S, Hajage D, Schortgen F, et al. Initiation strategies for renal-replacement therapy in the intensive care unit (AKIKI). \href{https://doi.org/10.1056/NEJMoa1603017}{\textit{N Engl J Med}. 2016;375:122-133.}

    \item Barbar SD, Clere-Jehl R, Bourredjem A, et al. Timing of renal-replacement therapy in patients with acute kidney injury and sepsis (IDEAL-ICU). \href{https://doi.org/10.1056/NEJMoa1803213}{\textit{N Engl J Med}. 2018;379:1431-1442.}

    \item STARRT-AKI Investigators. Timing of initiation of renal-replacement therapy in acute kidney injury. \href{https://doi.org/10.1056/NEJMoa1914579}{\textit{N Engl J Med}. 2020;383:240-251.}

    \item Jaber S, Paugam C, Futier E, et al. Sodium bicarbonate therapy for patients with severe metabolic acidaemia in the intensive care unit (BICAR-ICU). \href{https://doi.org/10.1016/S0140-6736(18)31080-8}{\textit{Lancet}. 2018;392:31-40.}
\end{enumerate}

\end{document}
